% latexmk -pdf -pvc blatt12.tex
\documentclass{hott-übung}

\begin{document}

\setcounter{blattnummer}{12}
\setcounter{aufgabennummer}{0}

\blatt

\aufgabe{}
Sei \(X\) ein Typ und \(A \xleftarrow{f} C \xrightarrow{g} B\) ein Spann.
Konstruiere Äquivalenzen:
\begin{enumerate}
  \item \(\Coeq(\leer \rightrightarrows X)\simeq X\).
  \item \(\Coeq(1 \rightrightarrows 1)\simeq S^1\).
  \item \(\Coeq(C \rightrightarrows A\sqcup B) \simeq A \sqcup_C B\).
  {\tiny Was sind die parallelen Abbildungen?}
\end{enumerate}

\aufgabe{}
Im Folgenden sei \(A\) ein Typ und \(n \geq -2\).
Zeige:
\begin{enumerate}
  \item
    Genau dann ist \(A\) ein \((n+1)\)-zusammenhängender Typ, wenn \(\|A\|\) gilt und \(x=y\) für alle \(x,y : A\) ein \(n\)-zusammenhängender Typ ist.
  \item
    Sei \(A\) ein punktierter zusammenhängender Typ.
    Folgere, dass \(A\) genau dann \((n+1)\)-zusammenhängend ist, falls \(\Omega A\) ein \(n\)-zusammenhängender Typ ist.
\end{enumerate}  
Eine Abbildung \(f : B \to A\) heißt \emph{\(n\)-zusammenhängend}, wenn \(\fib_f(a)\) für jedes \(a : A\) ein \(n\)-zusammenhängender Typ ist.
Sei nun \(n \geq 0\).
\begin{enumerate}[resume]
  \item
    Sei \(f : B \to_\ast A\) eine punktierte Abbildung zwischen zusammenhängenden punktierten Typen.
    Folgere, dass \(f\) genau dann \((n+1)\)-zusammenhängend ist, wenn \(\Omega f\) eine \(n\)-zusammenhängende Abbildung ist.
    \tiny{Lange Fasersequenz!}
\end{enumerate}

% Sei \(A\) ein punktierter \((n+1)\)-zusammenhängender Typ.
% Zeige, dass \(\ast : \eins \to_* A\) eine \(n\)-zusammenhängende Abbildung ist.

\aufgabe{}
Folgere aus der Existenz der Hopf-Faserung \(\Omega^3 S^3\simeq \Omega^3 S^2 \).

\end{document}
